\section{Introduction}
\label{sec:introduction}

The B method~\cite{AbrialBBook} supports the development of software from
mathematical specifications through successive refinements. Its industrial
toolsets, notably Atelier~B~\cite{AtelierB}, generate proof obligations
expressing properties such as invariant preservation and refinement
correctness. Users discharge these obligations with automatic provers
and interactive proof. The quality of this interaction matters: a proof
backend should expose understandable statements, provide reusable
mathematics, and make the assumptions behind each proof step explicit.

Lean~4~\cite{lean} offers a complementary environment for this work. Its
kernel checks proof terms independently of the tactics that construct
them, and its metaprogramming framework supports domain-specific proof
interfaces. Mathlib~\cite{mathlib4} provides an extensive mathematical
library. Bringing these facilities to existing B developments requires
an interpretation of B's set-based expressions in Lean's typed logic and
a treatment of the partial operations that occur throughout B proofs.

Partiality is central to this translation. A function application requires
an argument in the function's domain; a cardinality requires a finite
set; a minimum requires a set with a least element. These conditions
depend on the logical context of the expression. An operation guard may
justify a partial application that would be meaningless elsewhere.
Moreover, an interactive proof can introduce additional terms whose
well-definedness must also be justified. Generating a separate collection
of side conditions is useful only if the subsequent reasoning respects
their connection to the terms they justify.

This paper presents \barrel (B Automated tRanslation for Reasoning
in Lean), a backend that connects Atelier~B proof obligations to Lean's
interactive proof environment. Users retain their B source files.
\barrel invokes Atelier~B when generation is required, reads the exported
obligations, and constructs propositions over a library of B operators.
Sets and relations remain available as mathematical objects rather than
being replaced by an opaque verification condition language. Users can
therefore combine B-specific lemmas with ordinary Lean tactics and Mathlib.

For its proof-carrying operators, \barrel represents well-definedness as
an argument of the operator itself. During translation, a missing argument
becomes an explicit proof obligation closed over the surrounding variables
and hypotheses. Once proved, that evidence is used in the translated
expression. Repeated conditions can share a theorem, and a collection of
extensible tactics attempts to discharge routine conditions automatically.
Each remaining principal or WD obligation can be proved separately, with
earlier facts available as lemmas. A final command checks completion and
publishes the resulting declarations.

We use a fresh sensor-monitoring machine to make the intended benefit
concrete. A zone is authorized only when sufficiently many sensors have
reported values and their readings lie within configured bounds. The
property combines partial maps, cardinalities, and extrema. Routine
well-definedness reasoning establishes that these expressions have a
meaning; the principal proofs establish that updates preserve the
meaningful authorization property. In particular, a batch of changed
readings must invalidate affected zones while leaving unaffected zones'
authorizations justified. This separation gives a direct way to assess
whether the backend handles WD bookkeeping while leaving the developer
to address the substantive argument.

The paper presents the following contributions:
\begin{itemize}
  \item an integrated Atelier~B-to-Lean backend and mathematical library
    for conducting B proofs with ordinary Lean tactics;
  \item a contextual encoding of supported partial operators with explicit
    proof dependencies, shared WD conditions, and extensible automation;
  \item an interactive workflow for proving and reusing individual
    obligations, with an explicit completion and publication boundary.
\end{itemize}
The case study evaluates the resulting framework through a complete
development. Its 80 source obligations generate 754 WD proof requests,
which sharing reduces to 66 retained conditions. With the case study's
reusable lemmas, 51 of these conditions are discharged automatically;
the remaining obligations are proved through the individual proof
commands. A transitive axiom audit checks every published declaration.
Industrial import experience is supplementary to this focused account
of proof construction and completion.

The account distinguishes the properties checked by Lean from the
connection to the source B development. Kernel checking establishes the
translated propositions under their assumptions; it does not by itself
verify the external obligation generator or the translation. We state
this boundary alongside the implemented operator coverage and the proof
completion discipline.
